\chapter{Morphismes \'etales}
\label{I}



Pour simplifier l'exposition, on suppose que tous les
pr\'esch\'emas envisag\'es sont localement noeth\'eriens, du
moins apr\`es le num\'ero~\ref{I.2}.

\section{Notions de calcul diff\'erentiel}
\label{I.1}
Soit $X$ un pr\'esch\'ema sur~$Y$, soit $\Delta_{X/Y}$ ou $\Delta$
\label{indnot:ab}le morphisme diagonal $X\to X\times_Y X$. C'est un morphisme
d'immersion, donc un morphisme d'immersion ferm\'ee de~$X$ dans un
ouvert $V$ de~$X\times_Y X$. Soit $\cal{I}_X$ l'id\'eal du
sous-pr\'esch\'ema ferm\'e correspondant \`a la diagonale
dans~$V$ (N.B. si on veut faire les choses intrins\`equement, sans
supposer $X$ s\'epar\'e sur~$Y$ --- hypoth\`ese qui serait
canularesque --- on devrait consid\'erer l'image inverse ensembliste
de~$\cal{O}_{X\times X}$ dans~$X$, et d\'esigner par $\cal{I}_X$
l'id\'eal d'augmentation dans ce dernier...). Le faisceau
$\cal{I}_X/\cal{I}_X^2$ peut \^etre regard\'e comme un faisceau
quasi-coh\'erent sur~$X$, on le d\'enote
par~$\mathit{\Omega}^1_{X/Y}$.
\label{indnot:ac}Il est de type fini si $X\to Y$ est de type fini. Il se comporte bien
par rapport \`a extension de la base $Y'\to Y$. On introduit aussi
les faisceaux $\cal{O}_{X\times_Y X}/\cal{I}_X^{n+1}=\cal{P}^n_{X/Y}$,
\label{indnot:ad}ce sont des faisceaux d'\emph{anneaux} sur~$X$, faisant de~$X$ un
pr\'esch\'ema qu'on peut d\'enoter par $\Delta_{X/Y}^n$
\label{indnot:ae}et appeler le \emph{\nieme voisinage infinitésimal de}~$X/Y$.
\index{voisinage infinit\'esimal de~$X/Y$|hyperpage}Le sorite en est d'une trivialit\'e totale, bien qu'il soit assez
long\footnote{\cf EGA IV~16.3.}; il serait prudent de n'en parler
qu'au moment o\`u on en dit quelque chose de serviable, avec les
morphismes lisses.

\section{Morphismes quasi-finis}
\label{I.2}
\begin{proposition}
\label{I.2.1}
Soit $A\to B$ un homomorphisme local (N.B. les anneaux sont maintenant
noeth\'eriens), $\goth{m}$ l'id\'eal maximal de~$A$. Conditions
\'equivalentes:
\begin{enumerate}
\item[(i)] $B/\goth{m}B$ est de dimension finie sur~$k=A/\goth{m}$.
\item[(ii)] $\goth{m}B$ est un id\'eal de d\'efinition, et $B/\goth{r}(B)=\kres(B)$ est une extension de~$k=\kres(A)$
\item[(iii)] Le compl\'et\'e $\widehat{B}$ est fini sur celui
$\widehat{A}$ de~$A$.
\end{enumerate}
\end{proposition}
On dit alors
que $B$ \emph{est quasi-fini}
\index{quasi-fini (morphisme, alg\`ebre)|hyperpage}sur~$A$. Un morphisme $f\colon X\to Y$
est dit quasi-fini en~$x$ (ou le $Y$-pr\'esch\'ema $f$ est dit
quasi-fini en~$x$) si $\cal{O}_x$ est quasi-fini
sur~$\cal{O}_{f(x)}$. Cela revient aussi \`a dire que $x$ est
\emph{isolé dans sa fibre}~$f^{-1}(x)$. Un morphisme est dit
quasi-fini s'il l'est en tout point\footnote{Dans EGA~II~6.2.3 on
suppose de plus $f$ de type fini.}.
\begin{corollaire}
\label{I.2.2}
Si $A$ est complet, quasi-fini \'equivaut \`a fini.

\end{corollaire}

On pourrait donner le sorite \textup{(i), (ii), (iii), (iv), (v),} habituel pour les morphismes quasi-finis, mais ce ne semble pas indispensable ici.

\section{Morphismes non ramifi\'es ou nets}
\label{I.3}
\begin{proposition}
\label{I.3.1}
Soit $f\colon X\to Y$ un morphisme de type fini, $x\in X$,
$y=f(x)$. Conditions \'equivalentes:
\begin{enumerate}
\item[(i)] $\cal{O}_x/\goth{m}_y\cal{O}_x$ est une extension finie
s\'eparable de~$\kres(y)$.
\item[(ii)] $\mathit{\Omega}^1_{X/Y}$ est nul en~$x$.
\item[(iii)] Le morphisme diagonal $\Delta_{X/Y}$ est une immersion
ouverte au voisinage de~$x$.
\end{enumerate}
\end{proposition}
Pour l'implication (i)$\To$(ii), on est ramen\'e
aussit\^ot par Nakayama au cas o\`u $Y=\Spec(k)$,
$X=\Spec(k')$, o\`u c'est bien connu et d'ailleurs trivial
sur la d\'efinition de s\'eparable; (ii)$\To$(iii)
d'apr\`es une caract\'erisation agr\'eable et facile des
immersions ouvertes, utilisant Krull; (iii)$\To$(i) car on est
encore ramen\'e au cas o\`u $Y=\Spec(k)$ et o\`u le
morphisme diagonal est une immersion ouverte partout. Il faut alors
prouver que $X$ est fini d'anneau s\'eparable sur~$k$, on est
ramen\'e pour ceci au cas o\`u $k$ est alg\'ebriquement
clos. Mais alors tout point ferm\'e de~$X$ est isol\'e (car
identique \`a l'image inverse de la diagonale par le morphisme $X\to
X\times_k X$ d\'efini par~$x$), d'o\`u le fait que $X$ est
fini. On peut supposer alors $X$ r\'eduit \`a un point,
d'anneau~$A$, donc $A\otimes_k A\to A$ est un isomorphisme, d'o\`u
$A=k$ cqfd.

\begin{definition}
\label{I.3.2}
\begin{enumerate}
\item[a)] On dit alors que $f$ est \emph{net},
\index{net (morphisme, alg\`ebre)|hyperpage}ou encore \emph{non ramifié},
\index{non ramifi\'e (morphisme, alg\`ebre)|hyperpage}en~$x$, ou que $X$ est net, ou encore non ramifi\'e, en $x$ sur~$Y$.
\item[b)] Soit $A\to B$ un homomorphisme local, on dit qu'il est
\emph{net}, ou \emph{non ramifié}, ou que $B$ est une alg\`ebre
locale \emph{nette}, ou \emph{non ramifiée} sur~$A$, si 
$B/\goth{r}(A)B$ est une extension finie s\'eparable de~$A/\goth{r}(A)$, \ie si $\goth{r}(A)B=\goth{r}(B)$
et $\kres(B)$ est une extension s\'eparable
de~$\kres(A)$\footnote{\Cf remords dans III~\ref{III.1.2}.}.
\end{enumerate}
\end{definition}

\begin{remarquesstar}
Le fait
que $B$ soit net sur~$A$ se reconna\^it d\'ej\`a sur
les compl\'et\'es de~$A$ et de~$B$. Net implique quasi-fini.
\end{remarquesstar}

\begin{corollaire}
\label{I.3.3}
L'ensemble des points o\`u $f$ est net est ouvert.
\end{corollaire}

\begin{corollaire}
\label{I.3.4}
Soient $X'$, $X$ deux pr\'esch\'emas de type fini sur~$Y$, et
$g\colon X'\to X$ un $Y$-morphisme. Si $X$ est net sur~$Y$, le
morphisme graphe $\Gamma_g\colon X'\to X\times_Y X$ est une immersion
ouverte.
\end{corollaire}

En effet, c'est l'image inverse du morphisme diagonal $X\to X\times_Y
X$ par
$$
g\times_Y\id_{X'}\colon X'\times_Y X\to X\times_Y X.
$$

On peut aussi introduire l'id\'eal annulateur $\goth{d}_{X/Y}$
\label{indnot:af}de~$\mathit{\Omega}^1_{X/Y}$, appel\'e id\'eal diff\'erente
\index{diff\'erente (id\'eal)|hyperpage}\index{ideal differente@id\'eal diff\'erente|hyperpage}de~$X/Y$; il d\'efinit un sous-pr\'esch\'ema ferm\'e de~$X$
qui, ensemblistement, est l'ensemble des points o\`u $X/Y$ est
ramifi\'e, \ie non net.

\begin{proposition}
\label{I.3.5}
\begin{enumerate}
\item[(i)] Une immersion est nette.
\item[(ii)] Le compos\'e de deux morphismes nets l'est.
\item[(iii)] Extension de base dans un morphisme net en est un autre.
\end{enumerate}
\end{proposition}

Se voit indiff\'eremment sur (ii) ou (iii) (le deuxi\`eme me
semble plus amusant). On peut bien entendu aussi pr\'eciser, en
donnant des \'enonc\'es ponctuels; ce n'est plus g\'en\'eral
qu'en apparence (sauf dans le cadre de la d\'efinition~b)), et
barbant. On obtient comme d'habitude des corollaires:

\begin{corollaires}
\label{I.3.6}
\begin{enumerate}
\item[(iv)] Produit cart\'esien de deux morphismes nets en est un
autre.
\item[(v)] Si $gf$ est net, $f$ est net.
\item[(vi)] Si $f$ est net, $f_{\text{\textup{réd}}}$ est net.
\end{enumerate}
\end{corollaires}

\begin{proposition}
\label{I.3.7}
Soit $A\to B$ un homomorphisme local, on suppose l'extension
r\'esiduelle $\kres(B)/\kres(A)$ triviale ou $\kres(A)$ alg\'ebriquement
clos. Pour que $B/A$ soit net, il faut et il suffit que $\widehat{B}$
soit (comme $\widehat{A}$-alg\`ebre) un quotient de~$\widehat{A}$.
\end{proposition}

\begin{remarquesstar}
\begin{enumerate}
\item[--] Dans le cas o\`u on ne suppose pas l'extension
r\'esiduelle triviale, on peut se ramener \`a ce cas en faisant
une extension finie plate convenable sur~$A$ qui d\'etruise ladite
extension.
\item[--] Donner l'exemple o\`u $A$ est l'anneau local d'un point
double ordinaire d'une courbe, $B$ d'un point du normalis\'e: alors
$A\subset B$, $B$ est net sur
$A$ \`a extension r\'esiduelle triviale, et
$\widehat{A}\to\widehat{B}$ est surjectif mais \emph{non injectif}. On
va donc renforcer la notion de nettet\'e.
\end{enumerate}
\end{remarquesstar}

\section{Morphismes \'etales. Rev\^etements \'etales}
\label{I.4}
On va admettre tout ce qui nous sera n\'ecessaire sur les morphismes
plats; ces faits seront d\'emontr\'es ult\'erieurement, s'il y a
lieu\footnote{\Cf Exp.\, \ref{IV}.}.

\begin{definition}
\label{I.4.1}
\begin{enumerate}
\item[a)] Soit $f\colon X\to Y$ un morphisme de type fini. On dit que
$f$ est \emph{étale}
\index{etale (morphisme, algebre)@\'etale (morphisme, alg\`ebre)|hyperpage}en $x$ si $f$ est plat en $x$ et net en~$x$. On dit que $f$ est
\'etale s'il l'est en tous les points. On dit que $X$ est \'etale
en $x$ sur~$Y$, ou que c'est un $Y$-pr\'esch\'ema \'etale en~$x$
etc.
\item[b)] Soit $f\colon A\to B$ un homomorphisme local. On dit que $f$
est \'etale, ou que $B$ est \'etale sur~$A$, si $B$ est plat et
non ramifi\'e sur~$A$.\footnote{\Cf remords dans III~\ref{III.1.2}.}
\end{enumerate}
\end{definition}

\begin{proposition}
\label{I.4.2}
Pour que $B/A$ soit \'etale, il faut et il suffit que
$\widehat{B}/\widehat{A}$ le soit.
\end{proposition}

En
effet, c'est vrai s\'epar\'ement pour \og net\fg et pour \og plat\fg.

\begin{corollaire}
\label{I.4.3}
Soit $f\colon X\to Y$ de type fini, et $x\in X$. Le fait que $f$ soit
\'etale en $x$ ne d\'epend que de l'homomorphisme local
$\cal{O}_{f(x)}\to \cal{O}_x$, et m\^eme seulement de
l'homomorphisme correspondant pour les compl\'et\'es.
\end{corollaire}

\begin{corollaire}
\label{I.4.4}
Supposons que l'extension r\'esiduelle $\kres(A)\to \kres(B)$ soit triviale,
ou que $\kres(A)$ soit alg\'ebriquement clos. Alors $B/A$ est \'etale
\sss $\widehat{A}\to\widehat{B}$ est un isomorphisme.
\end{corollaire}

On conjugue la platitude et~\ref{I.3.7}.
\begin{proposition}
\label{I.4.5}
Soit $f\colon X\to Y$ un morphisme de type fini. Alors l'ensemble des
points o\`u il est \'etale est ouvert.
\end{proposition}

En effet, c'est vrai s\'epar\'ement pour \og net\fg et \og plat\fg.

Cette proposition montre qu'on peut laisser tomber les \'enonc\'es
\og ponctuels\fg dans l'\'etude des morphismes de type fini \'etales
quelque part.
\begin{proposition}
\label{I.4.6}
\begin{enumerate}
\item[(i)] Une immersion ouverte est \'etale.
\item[(ii)] Le compos\'e de deux morphismes \'etales est
\'etale.
\item[(iii)] Extension de la base.
\end{enumerate}
\end{proposition}

En effet, (i) est trivial, et pour (ii) et (iii) il suffit de noter
que c'est vrai pour \og net\fg et pour \og plat\fg. \`A vrai dire, il y a
aussi des \'enonc\'es correspondants pour les homomorphismes
locaux (sans condition de finitude), qui en tout \'etat de cause
devront figurer au multiplodoque (\`a commencer par le cas: net).

\begin{corollaire}
\label{I.4.7}
Un produit cart\'esien de deux morphismes \'etales est itou.
\end{corollaire}

\begin{corollaire}
\label{I.4.8}
Soient $X$ et $X'$ de type fini sur~$Y$, $g\colon X\to X'$ un
$Y$-morphisme. Si $X'$ est non ramifi\'e sur~$Y$ et $X$ \'etale
sur~$Y$, alors $g$ est \'etale.
\end{corollaire}

En effet, $g$ est le compos\'e du morphisme graphe $\Gamma_g\colon
X\to X\times_Y X'$ qui est une immersion ouverte par~\ref{I.3.4}, et
du morphisme de projection qui est \'etale car d\'eduit du
morphisme \'etale $X\to Y$ par le \og changement de base\fg $X'\to Y$.

\begin{definition}
\label{I.4.9}
On appelle rev\^etement \'etale (\resp net)
\index{etale (revetement)@\'etale (rev\^etement)|hyperpage}\index{net (revetement)@net (rev\^etement)|hyperpage}\index{revetement etale (\resp net)@rev\^etement \'etale (\resp net)|hyperpage}de~$Y$ un $Y$-sch\'ema $X$ qui est fini sur~$Y$ et \'etale
(\resp net) sur~$Y$.
\end{definition}

\enlargethispage{.5cm}La premi\`ere condition signifie que $X$ est d\'efini par un
faisceau coh\'erent d'alg\`ebres $\cal{B}$ sur~$Y$. La
deuxi\`eme signifie alors que $\cal{B}$ est localement libre sur~$Y$
(\resp rien du tout), \emph{et} que de plus, pour tout $y\in Y$, la
fibre $\cal{B}(y)=\cal{B}_y\otimes_{\cal{O}_y}\kres(y)$ soit une
alg\`ebre s\'eparable (= compos\'e fini d'extensions finies
s\'eparables) sur~$\kres(y)$.

\begin{proposition}
\label{I.4.10}
Soit $X$ un rev\^etement plat de~$Y$ de degr\'e $n$ (la
d\'efinition de ce terme m\'eritait de figurer dans~\ref{I.4.9})
d\'efini par un faisceau coh\'erent localement libre $\cal{B}$
d'alg\`ebres. On d\'efinit de fa\c con bien connue
l'homomorphisme trace $\cal{B}\to\cal{A}$ (qui est un homomorphisme de
$\cal{A}$-modules, o\`u $\cal{A}=\cal{O}_Y$). Pour que $X$ soit
\'etale, il \fets que la forme bilin\'eaire
$\trace_{\cal{B}/\cal{A}}xy$ correspondante d\'efinisse un
isomorphisme de~$\cal{B}$ sur~$\cal{B}$, ou ce qui revient au
m\^eme, que la \emph{section discriminant}
$$
d_{X/Y} = d_{\cal{B}/\cal{A}} \in
\Gamma(Y,\tbigwedge^n\check{\cal{B}}\otimes_{\cal{A}}
\tbigwedge^n\check{\cal{B}})
$$
soit inversible, ou enfin que l'id\'eal discriminant d\'efini par
cette section soit l'id\'eal unit\'e.
\end{proposition}
En effet, on est ramen\'e au cas o\`u $Y=\Spec(k)$, et alors c'est
un crit\`ere de s\'eparabilit\'e bien connue, et trivial par
passage \`a la cl\^oture alg\'ebrique de~$k$.

\begin{remarquestar}
On
aura un \'enonc\'e moins trivial plus bas, quand on ne suppose
pas a priori que $X$ est plat sur~$Y$ mais qu'on fait une
hypoth\`ese de normalit\'e.
\end{remarquestar}

\section{La propri\'et\'e fondamentale des morphismes \'etales}
\label{I.5}

\begin{theoreme}
\label{I.5.1}
Soit $f\colon X \to Y$ un morphisme de type fini. Pour que $f$ soit
une immersion ouverte, il \fets que ce soit un
morphisme \emph{étale et radiciel}.
\end{theoreme}

Rappelons que radiciel
\index{radiciel (morphisme)|hyperpage}signifie: injectif, \`a extensions
r\'esiduelles radicielles (et
on
peut aussi rappeler que cela signifie que le morphisme reste injectif
par toute extension de la base). La n\'ecessit\'e est triviale,
reste la suffisance. On va donner deux d\'emonstrations
diff\'erentes, la premi\`ere plus courte, la deuxi\`eme plus
\'el\'ementaire.

1) Un morphisme plat est ouvert, donc on peut supposer (rempla\c cant $Y$ par~$f(X)$) que $f$ est un \emph{homéomorphisme}
sur.
Par toute extension de base, il restera vrai que
$f$
est plat, radiciel, surjectif, donc un hom\'eomorphisme, a fortiori
ferm\'e
Donc $f$ est \emph{propre}. Donc $f$ est \emph{fini}
(r\'ef\'erence: th\'eor\`eme de Chevalley) d\'efini par un
faisceau coh\'erent $\mathcal{B}$ d'alg\`ebres. $\mathcal{B}$ est
localement libre, de plus en vertu de l'hypoth\`ese il est partout
de rang~$1$, donc $X=Y$, cqfd.

2) On peut supposer $Y$ et $X$ \emph{affines}. On se ram\`ene de
plus facilement \`a prouver ceci: si $Y = \Spec(A)$, $A$ local, et
si $f^{-1}(y)$ est non vide ($y$ \'etant le point ferm\'e de~$Y$)
alors $X=Y$ (en effet, cela impliquera que tout $y \in f(X)$ a un
voisinage ouvert $U$ tel que $X|U = U$). On aura $X = \Spec(B)$, on
veut prouver $A=B$. Mais pour ceci, on est ramen\'e \`a prouver
l'assertion analogue en rempla\c cant $A$ par~$\hat A$, $B$ par $B
\otimes_A \hat A$ (compte tenu que $\hat A$ est fid\`element plat
sur~$A$). On peut donc supposer $A$ \emph{complet}. Soit $x$ le point
au-dessus de~$y$, d'apr\`es le corollaire~\ref{I.2.2}
$\mathcal{O}_x$ est fini sur~$A$ donc (\'etant plat et radiciel
sur~$A$) est identique \`a~$A$. Donc on a $X = Y \amalg X^\prime$
(somme disjointe). Comme $X$ est radiciel sur~$Y$, $X^\prime$ est
vide. On a fini.
\enlargethispage{.5cm}\begin{corollaire}
\label{I.5.2}
Soit $f\colon X \to Y$ un morphisme d'\emph{immersion fermée} et
\emph{étale}. Si $X$ est connexe, $f$ est un isomorphisme de~$X$
sur une composante connexe de~$Y$.
\end{corollaire}

En effet, $f$ est aussi une immersion ouverte. On en d\'eduit:

\begin{corollaire}
\label{I.5.3}
Soit
$X$ un $Y$-sch\'ema net, $Y$ connexe. Alors toute section
de~$X$ sur~$Y$ est un isomorphisme de~$Y$ sur une composante connexe
de~$X$. Il y a donc correspondance biunivoque entre l'ensemble de ces
sections, et l'ensemble des composantes connexes $X_i$ de~$X$ telles
que la projection $X_i \to Y$ soit un isomorphisme, (ou, ce qui
revient au m\^eme par~\ref{I.5.1}, surjectif et radiciel). En
particulier, une section est connue quand en conna\^it sa valeur en
un point.
\end{corollaire}

Seule la premi\`ere assertion demande une d\'emonstration;
d'apr\`es~\ref{I.5.2} il suffit de remarquer qu'une section est une
immersion ferm\'ee (car $X$ est s\'epar\'e sur~$Y$) et \'etale
en vertu de~\ref{I.4.8}.

\begin{corollaire}
\label{I.5.4}
Soient $X$ et $Y$ deux pr\'esch\'emas sur~$S$, $X$ net
s\'epar\'e sur~$S$ et $Y$ connexe. Soient $f$, $g$ deux
$S$-morphismes de~$Y$ dans~$X$, $y$ un point de~$Y$, on suppose $f(y)
= g(y) = x$ et les homomorphismes r\'esiduels $\kres(x) \to \kres(y)$
d\'efinis par $f$ et~$g$ identiques (\og $f$ et~$g$ co\"incident
g\'eom\'etriquement en~$y$\fg). Alors $f$ et~$g$ sont identiques.
\end{corollaire}

R\'esulte de~\ref{I.5.3} en se ramenant au cas o\`u $Y=S$, en
rempla\c cant $X$ par $X \times_S Y$.

Voici une variante particuli\`erement importante de~\ref{I.5.3}.

\begin{theoreme}
\label{I.5.5}
Soient $S$ un pr\'esch\'ema, $X$ et $Y$ deux
$S$-pr\'esch\'emas, $S_0$ un sous-pr\'esch\'ema ferm\'e
de~$S$ ayant m\^eme espace sous-jacent que~$S$, $X_0 = X \times_S
S_0$ et $Y_0 = Y \times_S S_0$ les \og restrictions\fg de~$X$ et~$Y$
sur~$S_0$. On suppose $X$ \'etale sur~$S$. Alors l'application
naturelle
$$
\Hom_S(Y,X) \to \Hom_{S_0}(Y_0,X_0)
$$
est \emph{bijective}.
\end{theoreme}

On est encore ramen\'e au cas o\`u $Y=S$, et alors cela
r\'esulte de la description \og topologique\fg des sections de~$X/Y$
donn\'ee dans~\ref{I.5.3}.

\begin{scholiestar}
Ce r\'esultat comporte une assertion d'\emph{unicité} et
d'\emph{existence} de \emph{morphismes}. Il peut aussi s'exprimer
(lorsque $X$ et~$Y$ sont tous deux pris \'etales sur~$S$) que le
foncteur $X \mto X_0$ de la cat\'egorie des $S$-sch\'emas
\'etales dans la cat\'egorie des $S_0$-sch\'emas \'etales est
\emph{pleinement fidèle}, \ie \'etablit une
\emph{équivalence} de la premi\`ere avec une
\emph{sous-catégorie pleine} de la seconde. Nous verrons plus bas
que c'est m\^eme une \'equivalence de la premi\`ere et de la
seconde (ce qui sera un th\'eor\`eme d'\emph{existence de
$S$-schémas étales}).
\end{scholiestar}

La
forme suivante, plus g\'en\'erale en apparence, de~\ref{I.5.5}.\
est souvent commode:

\begin{corollaire}[\og Th\'eor\`eme de prolongement des rel\`evements\fg]
\label{I.5.6}
\index{prolongement des rel\`evements (th\'eor\`eme de)|hyperpage}\index{theoreme de prolongement des relevements@th\'eor\`eme de prolongement des rel\`evements|hyperpage}Consid\'erons un diagramme commutatif
$$
\xymatrix{{X}\ar[d]&{Y_0}\ar[l]\ar[d]\\{S}&{Y}\ar[l]}
$$
de morphismes, o\`u $X \to S$ est \'etale et $Y_0 \to Y$ est une
immersion ferm\'ee bijective. Alors on peut trouver un morphisme
unique $Y \to X$ qui rende les deux triangles correspondants
commutatifs.
\end{corollaire}

En effet, rempla\c cant $S$ par~$Y$ et $X$ par $X \times_S Y$, on
est ramen\'e au cas o\`u $Y=S$, et alors c'est le cas particulier
de~\ref{I.5.5} pour $Y=S$.

Signalons aussi la cons\'equence imm\'ediate suivante
de~\ref{I.5.1} (que nous n'avons pas
donn\'ee
en corollaire~1 pour ne pas interrompre la ligne d'id\'ees
d\'evelopp\'ee \`a la suite de~\ref{I.5.1}):

\begin{proposition}
\label{I.5.7}
Soient $X$, $X^\prime$ deux pr\'esch\'emas de type finis et plats
sur~$Y$, et soit $g\colon X \to X^\prime$ un $Y$-morphisme. Pour que
$g$ soit une immersion ouverte (\resp un isomorphisme) il faut et il
suffit que pour tout $y\in Y$, le morphisme induit sur les fibres
$$
g \otimes_Y \kres(y)\colon X \otimes_Y \kres(y) \to X^\prime
\otimes_Y \kres(y)
$$
le soit.
\end{proposition}

Il suffit de prouver la suffisance; comme c'est vrai pour la notion de
surjection, on est ramen\'e au cas d'une immersion
ouverte. D'apr\`es~\ref{I.5.1}, il faut v\'erifier que $g$ est
\emph{radiciel}, ce qui est trivial, et qu'il est \emph{étale}, ce
qui r\'esulte du corollaire~\ref{I.5.9} ci-dessous.

\begin{corollaire}
\label{I.5.8}
(devrait passer au \No\ref{I.3}) Soient $X$ et $X^\prime$ deux
$Y$-pr\'esch\'emas, $g\colon X \to X^\prime$ un $Y$-morphisme, $x$
un point de~$X$ et $y$ sa projection sur~$Y$. Pour que $g$ soit
quasi-fini (\resp net) en~$x$, il \fets qu'il en soit
de m\^eme de~$g\otimes_Y \kres(y)$.
\end{corollaire}

En effet, les deux alg\`ebres sur~$k\big(g(x)\big)$ qu'il faut
regarder pour s'assurer que l'on a bien un morphisme quasi-fini \resp net en~$x$ sont les m\^emes pour~$g$ et $g \otimes_Y \kres(y)$.

\begin{corollaire}
\label{I.5.9}
Avec
les notations de~\ref{I.5.8}, supposons $X$ et~$X^\prime$ plats
et de type fini sur~$Y$. Pour que $g$ soit plat (\resp \'etale)
en~$x$, il \fets que $g \otimes_Y \kres(y)$ le soit.
\end{corollaire}

Pour \og plat\fg l'\'enonc\'e n'est mis que pour m\'emoire, c'est
un des crit\`eres fondamentaux de platitude\footnote{\Cf IV~\ref{IV.5.9}.}. Pour \'etale, cela en r\'esulte; compte tenu
de~\ref{I.5.8}.

\section{Application aux extensions \'etales des anneaux locaux complets}
\label{I.6}

Ce num\'ero est un cas particulier de r\'esultats sur les
pr\'esch\'emas formels, qui devront figurer dans le
multiplodoque. N\'eanmoins, on s'en tire ici \`a meilleur compte,
\ie sans la d\'etermination locale explicite des morphismes
\'etales au \No\ref{I.7} (utilisant le Main Theorem). C'est peut-\^etre
une raison suffisante de garder le pr\'esent num\'ero (m\^eme
dans le multiplodoque) \`a cette place.

\begin{theoreme}
\label{I.6.1}
Soit $A$ un anneau local complet (noeth\'erien bien s\^ur), de
corps r\'esiduel~$k$. Pour toute $A$-alg\`ebre~$B$, soit $R(B) = B
\otimes_A k$ consid\'er\'e comme $k$-alg\`ebre, elle d\'epend
donc fonctoriellement de~$B$. Alors $R$ d\'efinit une
\emph{équivalence} de la cat\'egorie des $A$-alg\`ebres
\emph{finies et étales sur~$A$} avec la cat\'egorie des
alg\`ebres \emph{de rang fini séparables} sur~$k$.
\end{theoreme}

Tout d'abord, le foncteur en question est pleinement fid\`ele, comme
il r\'esulte du fait plus g\'en\'eral:

\begin{corollaire}
\label{I.6.2}
Soient $B$, $B^\prime$ deux $A$-alg\`ebres finies sur~$A$. Si $B$
est \'etale sur~$A$, alors l'application canonique
$$
\Hom_{\textup{$A$-alg}}(B,B^\prime) \to
\Hom_{\textup{$k$-alg}}\big(R(B),R(B^\prime)\big)
$$
est bijective.
\end{corollaire}

On est ramen\'e au cas o\`u $A$ est artinien (en rempla\c cant
$A$ par~$A/\goth{m}^n$), et alors c'est un cas particulier
de~\ref{I.5.5}.

Il reste \`a prouver que pour toute $k$-alg\`ebre finie et
s\'eparable (pourquoi ne pas dire: \'etale, c'est plus court) $L$,
il existe un $B$ \'etale sur~$A$ tel que $R(B)$ soit isomorphe
\`a~$L$. On peut supposer que $L$ est une extension s\'eparable
de~$k$, comme telle elle admet un g\'en\'erateur~$x$, \ie est
isomorphe \`a une alg\`ebre $k[t]/Fk[t]$ o\`u $F \in k[t]$ est
un polyn\^ome unitaire. On rel\`eve~$F$ en un polyn\^ome
unitaire $F_1$ dans $A[t]$, et on prend $B = A[t]/F_1A[t]$.

\section{Construction locale des morphismes non ramifi\'es et \'etales}
\label{I.7}

\begin{proposition}
\label{I.7.1}
Soient $A$ un anneau noeth\'erien, $B$ une alg\`ebre finie
sur~$A$, $u$ un g\'en\'erateur de~$B$ sur~$A$, $F \in A[t]$ tel
que $F(u) = 0$ (on ne suppose pas $F$ unitaire), $u^\prime =
F^\prime(u)$ (o\`u $F^\prime$ est le polyn\^ome d\'eriv\'e),
$\goth{q}$ un id\'eal premier de~$B$ ne contenant pas~$u^\prime$,
$\goth{p}$ sa trace sur~$A$. Alors $B_\goth{q}$ est net sur~$A_\goth{p}$.
\end{proposition}

En d'autres termes, posant $Y = \Spec(A)$, $X= \Spec(B)$,
$X_{u^\prime} = \Spec(B_{u^\prime})$, $X_{u^\prime}$ est non
ramifi\'e sur~$Y$. L'\'enonce r\'esulte du suivant, plus
pr\'ecis:

\begin{corollaire}
\label{I.7.2}
L'id\'eal diff\'erente de~$B/A$ contient $u^\prime B$, et lui est
\'egal si l'homo\-morphisme naturel $A[t]/FA[t] \to B$ (appliquant
$t$ dans~$u$) est un isomorphisme.
\end{corollaire}

Soit $J$ le noyau de l'homomorphisme $C = A[t] \to B$, ce noyau
contient $FA[t]$, et lui est \'egal dans le deuxi\`eme cas
envisag\'e dans~\ref{I.7.2}. Comme il est surjectif,
$\Omega^1_{B/A}$ s'identifie au quotient de~$\Omega^1_{C/A}$ par le
sous-module engendr\'e par $J\Omega^1_{C/A}$ et $d(J)$ (il aurait
fallu expliciter au \No\ref{I.1} la d\'efinition de l'homomorphisme~$d$, et
le calcul de~$\Omega^1$ pour une alg\`ebre de polyn\^omes).
Identifiant $\Omega^1_{C/A}$ \`a~$C$ gr\^ace \`a la base~$dt$,
on trouve $B/B\cdot J^\prime$ donc la diff\'erente est engendr\'ee
par l'ensemble~$J^\prime$ des images dans~$B$ des d\'eriv\'es des
$G \in J$, (et il suffit de prendre des~$G$ engendrant~$J$). Comme
$F\in J$, \resp $F$ est un g\'en\'erateur de~$J$, on a
fini. (N.B.\ On devrait mettre~\ref{I.7.2} en prop.\
et~\ref{I.7.1} en corollaire). On trouve:

\begin{corollaire}
\label{I.7.3}
Sous les conditions de~\ref{I.7.1}, supposant $F$ unitaire et que
$A[t]/FA[t] \to B$ est un isomorphisme, pour que $B_\goth{q}$ soit
\'etale sur~$A_\goth{p}$, il
faut et il suffit
que~$\goth{q}$ ne contienne
pas~$u^\prime$.
\end{corollaire}

En effet, comme $B$ est plat sur~$A$, \'etale \'equivaut \`a
net, et on peut appliquer~\ref{I.7.2}.

\begin{corollaire}
\label{I.7.4}
Sous les conditions de~\ref{I.7.3} pour que $B$ soit \'etale
sur~$A$ il
faut et il suffit que~$u^\prime$ soit inversible, ou encore que l'id\'eal engendr\'e
par $F$, $F^\prime$ dans $A[t]$ soit l'id\'eal unit\'e.
\end{corollaire}

Le dernier crit\`ere r\'esulte du premier et de Nakayama
(dans~$B$).

Un polyn\^ome unitaire $F\in A[t]$ ayant la propri\'et\'e
\'enonc\'ee dans le corollaire~\ref{I.7.4}
est dit \emph{polynôme séparable}
\index{separable (polynome)@s\'eparable (polyn\^ome)|hyperpage}(si $F$ n'est pas unitaire, il faudrait au moins exiger que le
coefficient de son terme dominant soit inversible; dans le cas o\`u
$A$ est un corps, on retrouve la d\'efinition usuelle).


\begin{corollaire}
\label{I.7.5}
Soit $B$ une alg\`ebre \emph{finie} sur l'anneau
\emph{local}~$A$. On suppose que $K(A)$ est infini ou que $B$ soit
local. Soit~$n$ le rang de~$L = B \otimes_A K(A)$ sur~$K(A) = k$. Pour
que $B$ soit net (\resp \'etale) sur~$A$, il faut et il suffit que
$B$ soit isomorphe \`a un quotient de (\resp isomorphe \`a)
$A[t]/FA[t]$, o\`u $F$ est un polyn\^ome unitaire s\'eparable,
qu'on peut supposer (\resp qui est n\'ecessairement) de
degr\'e~$n$.
\end{corollaire}

Il n'y a qu'\`a prouver la n\'ecessit\'e. Supposons $B$ net
sur~$A$, donc $L$ s\'eparable sur~$k$, il r\'esulte alors de
l'hypoth\`ese faite que $L/k$ admet un g\'en\'erateur~$\xi$,
donc les $\xi^i$ ($0 \leq i < n$) forment une base de~$L$
sur~$k$. Soit $u\in B$ relevant~$\xi$, alors par Nakayama les $u^i$
($0\leq i < n$) engendrent le $A$-module~$B$ (\resp en forment une
base), en particulier on peut trouver un polyn\^ome unitaire $F \in
A[t]$ tel que $F(u)= 0$, et $B$ sera isomorphe \`a un quotient de
(\resp isomorphe \`a) $A[t]/FA[t]$. Enfin, en vertu
de~\ref{I.7.4}.\ appliqu\'e \`a $L/k$, $F$ et~$F^\prime$
engendrent $A[t]$
modulo
$\goth{m} A[t]$, donc (d'apr\`es Nakayama dans $A[t]/FA[t]$) $F$
et~$F^\prime$ engendrent~$A[t]$, on a fini.

\begin{theoreme}
\label{I.7.6}
Soient $A$ un anneau local, $A \to \cal{O}$ un homomorphisme local tel
que~$\cal{O}$ soit isomorphe \`a une alg\`ebre localis\'ee d'une
alg\`ebre de type fini sur~$A$. Supposons
$\cal{O}$
\emph{net} sur~$A$. Alors on peut trouver une $A$-alg\`ebre $B$,
enti\`ere sur~$A$, un id\'eal maximal $\goth{n}$ de~$B$, un
g\'en\'erateur $u$ de~$B$ sur~$A$, un polyn\^ome unitaire $F\in
A[t]$, tels que $\goth{n} \not\ni F^\prime(u)$ et que $\cal{O}$ soit
isomorphe (comme $A$-alg\`ebre) \`a~$B_{\goth{n}}$. Si $\cal{O}$
est \'etale sur~$A$, on peut prendre $B=A[t]/FA[t]$.
\end{theoreme}

(Bien entendu, on a l\`a des conditions aussi suffisantes...)

Signalons d'abord les agr\'eables corollaires:

\begin{corollaire}
\label{I.7.7}
Pour que $\cal{O}$ soit net sur~$A$, il \fets que $\cal{O}$ soit
isomorphe au quotient d'une alg\`ebre analogue et \emph{étale}
sur~$A$.
\end{corollaire}

En effet, on prendra $\cal{O}^\prime = B^\prime_{\goth{n}^\prime}$,
o\`u $B^\prime = A[t]/FA[t]$ et o\`u $\goth{n}^\prime$ est l'image
inverse de~$\goth{n}$ dans~$B^\prime$.

\begin{corollaire}
\label{I.7.8}
Soit
$f\colon X \to Y$ un morphisme de type fini, $x\in X$. Pour que
$f$ soit net en~$x$, il \fets qu'il existe un voisinage ouvert
$U$ de~$x$ tel que $f|U$ se factorise en $U \to X^\prime \to Y$,
o\`u la premi\`ere fl\`eche est une immersion ferm\'ee et la
seconde un morphisme \'etale.
\end{corollaire}

C'est une simple traduction de~\ref{I.7.7}.

Montrons comment le jargon de~\ref{I.7.6} r\'esulte de
l'\'enonc\'e principal: en effet, il existe par~\ref{I.7.7} un
\'epimorphisme $\cal{O}^\prime \to \cal{O}$, o\`u $\cal{O}$ a les
propri\'et\'es voulues; mais comme $\cal{O}^\prime$ et~$\cal{O}$
sont \'etales sur~$A$, le morphisme $\cal{O}^\prime \to \cal{O}$ est
\'etale par~\ref{I.4.8} donc un isomorphisme.

\subsubsection*{D\'emonstration de~\ref{I.7.6}}
Elle reprend une d\'emonstration du s\'eminaire
Chevalley. D'apr\`es le \emph{Main Theorem} on aura $\cal{O} =
B_{\goth{n}}$, o\`u $B$ est une alg\`ebre finie sur~$A$ et
$\goth{n}$ en est un id\'eal maximal. Alors $B/\goth{n} =
K(\cal{O})$ est une extension s\'eparable donc monog\`ene de~$k$;
si $\goth{n}_i$ ($1 \leq i \leq r$) sont les id\'eaux maximaux
de~$B$ distincts de~$\goth{n}$, il existe donc un \'el\'ement $u$
de~$B$ qui appartient \`a tous les~$\goth{n}_i$, et dont l'image
dans~$B/\goth{n}$ en est un g\'en\'erateur. Or $B/\goth{n} =
B_\goth{n}/\goth{n} B_\goth{n} = B_\goth{n}/\goth{m} B_\goth{n}$
(o\`u $\goth{m}$ est l'id\'eal maximal de~$A$). Admettons un
instant le

\begin{lemme}
\label{I.7.9}
Soient $A$ un anneau local, $B$ une alg\`ebre finie sur~$A$,
$\goth{n}$ un id\'eal maximal de~$B$, $u$ un \'el\'ement de~$B$
dont l'image dans $B_\goth{n}/\goth{m} B_\goth{n}$ l'engendre comme
alg\`ebre sur~$k = A/\goth{m}$, et qui se trouve dans tous les
id\'eaux maximaux de~$B$ distincts de~$\goth{n}$. Soit $B^\prime =
B[u]$, $\goth{n}^\prime = \goth{n} B^\prime$. Alors l'homomorphisme
canonique $B^\prime_{\goth{n}^\prime} \to B_\goth{n}$ est un
isomorphisme.
\end{lemme}

\begin{lemme}
\label{I.7.10}
(aurait d\^u figurer en corollaire \`a~\ref{I.7.1}
avant~\ref{I.7.5} qu'il implique). Soit $B$ une alg\`ebre finie
sur~$A$ engendr\'ee par un \'el\'ement~$u$, soit $\goth{n}$ un
id\'eal maximal de~$B$ tel que $B_\goth{n}$ soit non ramifi\'e
sur~$A$. Alors il existe un polyn\^ome unitaire $F \in A[t]$ tel que
$F(u) = 0$ et $F^\prime(u) \notin \goth{n}$.
\end{lemme}

Soit en effet $n$ le rang de la $k$-alg\`ebre $L = B \otimes_A k$,
d'apr\`es Nakayama il existe un polyn\^ome unitaire de
degr\'e~$n$ dans $A[t]$, tel que $F(u) = 0$. Soit $f$ le
polyn\^ome d\'eduit de~$F$ par r\'eduction mod~$\goth{m}$, alors
$L$ est $k$-isomorphe \`a $k[t]/fk[t]$, donc par~\ref{I.7.3}
$f^\prime(\xi)$ n'est pas contenu dans l'id\'eal maximal de~$L$ qui
correspond \`a~$\goth{n}$ ($\xi$ d\'esignant l'image de~$t$
dans~$L$, \ie l'image de~$u$ dans~$L$). Comme $f^\prime(\xi)$ est
l'image de~$F^\prime(u)$, on a fini.

Le
th\'eor\`eme~\ref{I.7.6} r\'esulte maintenant de la conjonction
de~\ref{I.7.9} et~\ref{I.7.10}. Reste \`a
prouver~\ref{I.7.9}. Posons $S^\prime = B^\prime - \goth{n}^\prime$,
donc $B^\prime {S^\prime}^{-1} = B^\prime_{\goth{n}^\prime}$.
$$
\begin{array}{ccccc}
B&\longrightarrow&B{S^\prime}^{-1}&\longrightarrow&BS^{-1} =
B_\goth{n}\\ \Big\uparrow&&\Big\uparrow&&\\
B^\prime&\longrightarrow&B^\prime{S^\prime}^{-1}
\rlap{\hbox{$\displaystyle \ = B^\prime_{\goth{n}^\prime}$}}&&
\end{array}
$$
Soit de m\^eme $S=B-\goth{n}$, donc $BS^{-1} = B_\goth{n}$, on a
donc un homomorphisme naturel $B{S^\prime}^{-1} \to BS^{-1} =
B_\goth{n}$, prouvons que c'est un isomorphisme, \ie que les
\'el\'ements de~$S$ sont inversibles dans $B{S^\prime}^{-1}$,
\ie que tout id\'eal maximal $\goth{p}$ de ce dernier ne rencontre
pas~$S$, \ie induit~$\goth{n}$ sur~$B$. En effet, comme
$B{S^\prime}^{-1}$ est fini sur~$B^\prime{S^\prime}^{-1} =
B^\prime_{\goth{n}^\prime}$, $\goth{p}$ induit l'unique id\'eal
maximal $\goth{n}^\prime B_{\goth{n}^\prime}$
de~$B^\prime_{\goth{n}^\prime}$, donc induit l'id\'eal maximal
$\goth{n}^\prime$ de~$B^\prime$; comme $B$ est fini sur~$B^\prime$,
l'id\'eal $\goth{q}$ de~$B$ induit par~$\goth{p}$ \'etant
au-dessus de~$\goth{n}^\prime$, est n\'ecessairement maximal, et ne
contient pas~$u$, donc est identique \`a~$\goth{n}$. (On vient
d'utiliser que $u$ appartient \`a tout id\'eal maximal de~$B$
distinct de~$\goth{n}$). Prouvons maintenant que $B{S^\prime}^{-1}$
\'egale $B^\prime{S^\prime}^{-1}$: comme il est fini sur ce dernier,
on est ramen\'e par Nakayama \`a prouver l'\'egalit\'e mod
$\goth{n}^\prime B {S^\prime}^{-1}$ et a fortiori il suffit de prouver
l'\'egalit\'e mod $\goth{m} B{S^\prime}^{-1}$; or
$B{S^\prime}^{-1}/\goth{m} B{S^\prime}^{-1} = B_\goth{n}/\goth{m}
B_\goth{n}$ est engendr\'e sur~$k$ par~$u$ (on utilise ici l'autre
propri\'et\'e de~$u$) donc l'image de~$B^\prime$ (et a fortiori de
$B^\prime {S^\prime}^{-1}$) dedans est tout (comme sous-anneau
contenant~$k$ et l'image de~$u$).

\begin{remarquestar}
On doit pouvoir \'enoncer le th\'eor\`eme~\ref{I.7.6} pour un
anneau $\cal{O}$ qui est seulement semi-local, de fa\c con \`a
coiffer aussi~\ref{I.7.5}: on fera l'hypoth\`ese que
$\cal{O}/\goth{m}\cal{O}$ est une $k$-alg\`ebre \emph{monogène};
on pourra donc trouver un $u\in B$ dont l'image dans
$B/\goth{m}B$
est un g\'en\'erateur, et appartenant \`a tous les id\'eaux
maximaux de~$B$ ne provenant pas de~$\cal{O}$. Les lemmes~\ref{I.7.9}
et~\ref{I.7.10} doivent s'adapter sans difficult\'e. Plus
g\'en\'eralement, ...
\end{remarquestar}

\section[Rel\`evement infinit\'esimal des sch\'emas \'etales]{Rel\`evement infinit\'esimal des sch\'emas \'etales.
Application aux sch\'emas formels}
\label{I.8}

\begin{proposition}
\label{I.8.1}
Soient $Y$ un pr\'esch\'ema, $Y_0$ un sous-pr\'esch\'ema,
$X_0$ un $Y_0$-sch\'ema \'etale, $x$ un point de~$X_0$. Alors il
existe un $Y$-sch\'ema \'etale~$X$, un voisinage $U_0$ de~$x$
dans~$X_0$, et un $Y_0$-isomorphisme $U_0 \isomto X \times_Y Y_0$.
\end{proposition}

Soit en effet $y$ la projection de~$x$ dans~$Y_0$,
appliquant~\ref{I.7.6} au homomorphisme local \'etale $A_0 \to B_0$
des anneaux locaux de~$y$ et~$x$ dans $Y_0$ et~$X_0$: on trouve un
isomorphisme$$
B_0 = {C_0}_{\goth{n}_0}\qquad C_0 = A_0[t]/F_0A_0[t]
$$
o\`u $F_0$ est un polyn\^ome unitaire et $\goth{n}_0$ est un
id\'eal maximal de~$C_0$ ne contenant pas la classe de
$F^\prime_0(t)$ dans~$C_0$. Soit $A$ l'anneau local de~$y$ dans~$Y$,
soit $F$ un polyn\^ome unitaire dans $A[t]$ donnant~$F_0$ par
l'homomorphisme surjectif $A \to A_0$ (on rel\`eve les coefficients
de~$F_0$), soit enfin
$C = A[t]/FA[t]$
et $\goth{n}$ l'id\'eal maximal
de~$C$
image inverse de
$\goth{n}_0$
par l'\'epimorphisme naturel $C \to C \otimes_A A_0 = C_0$. Posons
$$
B = C_\goth{n}.
$$
Il est imm\'ediat par construction et~\ref{I.7.1} que $B$ est
\'etale sur~$A$, et qu'on a un isomorphisme $B \otimes_A A_0 =
A_0$. On sait (Chap.\, I)
qu'il existe un $Y$-sch\'ema de type fini~$X$ et un point~$z$ de~$X$
au-dessus de~$y$ tel que $\cal{O}_z$ soit $A$-isomorphe \`a~$C$;
comme ce dernier est \'etale sur~$A = \cal{O}_y$, on peut (en
prenant $X$ assez petit) supposer que $X$ est \'etale sur~$Y$. Soit
$X_0^\prime = X \times_Y Y_0$, alors l'anneau local de~$z$
dans~$X_0^\prime$ s'identifie \`a $\cal{O}_z \otimes_A A_0 = B
\otimes_A A_0$, donc est isomorphe \`a~$B_0$. Cet isomorphisme est
d\'efini pas un isomorphisme d'un voisinage~$U_0$ de~$x$ dans~$X_0$
sur un voisinage de~$z$ dans~$X_0^\prime$ (loc.\ cit\'e), qu'on peut
supposer identique \`a~$X_0^\prime$ en prenant $X$ assez petit. On a
fini.

\begin{corollaire}
\label{I.8.2}
\'Enonc\'e analogue pour des \emph{revêtements} \'etales, en
supposant le corps r\'esiduel $\kres(y)$ infini.
\end{corollaire}

La d\'emonstration est la m\^eme, \ref{I.7.5} rempla\c cant~\ref{I.7.6}.

\begin{theoreme}
\label{I.8.3}
Le foncteur envisag\'e dans~\ref{I.5.5} est une
\emph{équivalence} de \emph{catégories}.
\end{theoreme}

En vertu du th\'eor\`eme~\ref{I.5.5}, il reste \`a montrer que
tout $S_0$-sch\'ema \'etale~$X_0$ est isomorphe \`a un
$S_0$-sch\'ema $X \times_S S_0$, o\`u $X$ est un $S$-sch\'ema
\'etale. L'espace topologique sous-jacent \`a~$X$ devra \^etre
n\'ecessairement identique \`a celui de~$X_0$, $X_0$ s'identifiant
de plus \`a un sous-pr\'esch\'ema ferm\'e de~$X$. Le
probl\`eme est donc \'equivalent au suivant: trouver sur l'espace
topologique sous-jacent $|X_0|$ \`a~$X_0$ un faisceau
d'alg\`ebres~$\cal{O}_X$ sur~$f_0^*(\cal{O}_S)$ (o\`u $f_0$ est
la projection $X_0 \to S_0$, regard\'ee ici comme application
continue des espaces sous-jacents), faisant de~$|X_0|$ un
$S$-pr\'esch\'ema \'etale~$X$, et un homomorphisme
d'alg\`ebres $\cal{O}_X \to \cal{O}_{X_0}$, compatible avec
l'homomorphisme $f_0^*(\cal{O}_S) \to f_0^*(\cal{O}_{S_0})$ sur
les faisceaux de scalaires, induisant un isomorphisme $\cal{O}_X
\otimes_{f_0^*(\cal{O}_S)} f_0^*(\cal{O}_{S_0}) \isomto
\cal{O}_{X_0}$. (Alors $X$ sera un $S$-pr\'esch\'ema \'etale se
r\'eduisant suivant~$X_0$,
donc sera s\'epar\'e sur~$S$ puisque $X_0$ l'est sur~$S_0$, et $X$
r\'epond \`a la question). Si d'ailleurs $(U_i)$ est un
recouvrement de~$X_0$ par des ouverts, et si on a trouv\'e une
solution du probl\`eme dans chacun des~$U_i$, il r\'esulte du
th\'eor\`eme d'unicit\'e~\ref{I.5.5} que ces solutions se
recollent (\ie les faisceaux d'alg\`ebres qui les d\'efinissent,
munis de leurs homomorphismes d'augmentation, se recollent), et on
constate aussit\^ot que l'espace annel\'e ainsi construit
au-dessus de~$S$ est un $S$-pr\'esch\'ema \'etale~$X$ muni d'un
isomorphisme $X \times_S S_0 \isomfrom X_0$. Il suffit donc de trouver
une solution localement, ce qui est assur\'e par~\ref{I.8.1}.

\begin{corollaire}
\label{I.8.4}
Soient $S$ un pr\'esch\'ema formel localement noeth\'erien, muni
d'un id\'eal de d\'efinition~$J$, soit $S_0 =
\big(|S|,\cal{O}_S/\cal{J}\big)$ le pr\'esch\'ema ordinaire
correspondant. Alors le foncteur $\goth{X} \mto \goth{X} \times_S
S_0$ de la cat\'egorie des rev\^etements \'etales de~$S$ dans la
cat\'egorie des rev\^etements \'etales de~$S_0$ est une
\'equivalence de cat\'egories.
\end{corollaire}

Bien entendu, on appelera rev\^etement \'etale d'un
pr\'esch\'ema \emph{formel}~$S$ un rev\^etement de~$S$, \ie un
pr\'esch\'ema formel sur~$S$ d\'efini \`a l'aide d'un faisceau
coh\'erent d'alg\`ebres~$\cal{B}$, tel que $\cal{B}$ soit
\emph{localement libre} et que les fibres r\'esiduelles $\cal{B}_s
\otimes_{\cal{O}_s} \kres(s)$ de~$\cal{B}$ soient des alg\`ebres
\emph{séparables} sur~$\kres(s)$. Si on d\'esigne par~$S_n$ le
pr\'esch\'ema ordinaire
$\big(|S|,\cal{O}_S/J^{n+1}\big)$,
la donn\'ee d'un faisceau coh\'erent d'alg\`ebres $\cal{B}$
sur~$S$ \'equivaut \`a la donn\'ee d'une suite de faisceaux
coh\'erents d'alg\`ebres
$\cal{B}_n$
sur
les~$S_n$,
munis d'un syst\`eme transitif d'homomorphismes
$\cal{B}_m\to \cal{B}_n$
($m \geq n$) d\'efinissant des isomorphismes
$\cal{B}_m \otimes_{\cal{O}_{S_m}} \cal{O}_{S_n} \isomto \cal{B}_n$.
Il est imm\'ediat que $\cal{B}$ est localement libre si et seulement
si les $\cal{B}_n$ sur les~$S_n$ le sont, et que la condition de
s\'eparabilit\'e est v\'erifi\'ee si et seulement si elle
l'est pour~$\cal{B}_0$, ou encore pour tous
les~$\cal{B}_n$.
Ainsi, $\cal{B}$ est \'etale sur~$S$ si et seulement si les
$\cal{B}_n$
sur les~$S_n$ le sont. Compte tenu de cela, \ref{I.8.4} r\'esulte
aussit\^ot de~\ref{I.8.3}.

\begin{remarquestar}
Il n'\'etait pas n\'ecessaire dans~\ref{I.8.4} de se borner au cas
des \emph{revêtements}. C'est cependant le seul utilis\'e pour
l'instant.
\end{remarquestar}

\section{Propri\'et\'es de permanence}
\label{I.9}
Soit $A\to B$ un homomorphisme local et \'etale, nous examinons ici
quelques cas o\`u une certaine propri\'et\'e pour $A$
entra\^ine la m\^eme propri\'et\'e pour~$B$, ou
r\'eciproquement.
Un certain nombre de telles propositions sont d\'ej\`a
cons\'equences du simple fait que $B$ est \emph{quasi-fini} et
\emph{plat} sur~$A$, et nous nous bornerons \`a en \og rappeler\fg
quelques-unes: $A$~\emph{et} $B$ \emph{ont même dimension de
Krull, et même profondeur} (\og codimension cohomologique\fg de
Serre, dans la terminologie encore courante). Il en r\'esulte par
exemple que $A$ \emph{est Cohen-Macaulay si et seulement si $B$ l'est}.
D'ailleurs, pour tout id\'eal premier $\goth{q}$ de~$B$, induisant
$\goth{p}$ sur~$A$, $B_{\goth{q}}$ sera encore quasi-fini et plat
sur~$A_{\goth{p}}$, pourvu qu'on suppose que $B$ soit localis\'ee
d'une alg\`ebre de type fini sur~$A$ (cela r\'esulte du fait que
l'ensemble des points o\`u un morphisme de type fini est quasi-fini
\resp plat est ouvert); et d'ailleurs \emph{tout} id\'eal premier
$\goth{p}$ de~$A$ est induit par un id\'eal premier $\goth{q}$
de~$B$ (car $B$ est \emph{fidèlement} plat sur~$A$). Il en
r\'esulte par exemple que \emph{$\goth{p}$ et~$\goth{q}$ ont
même rang}; et encore que \emph{$A$ est sans idéal premier
immergé si et seulement si~$B$ l'est}.

Nous allons nous borner donc aux propositions plus sp\'eciales au
cas des morphismes \'etales.

\begin{proposition}
\label{I.9.1}
Soit $A\to B$ un homomorphisme local \'etale. Pour que $A$ soit
r\'egulier, il faut et il suffit que $B$ le soit.
\end{proposition}
En effet, soit $k$ le corps r\'esiduel de~$A$, $L$ celui de~$B$.
Comme $B$ est plat sur~$A$ et que
$L=B\otimes_A k$,
\ie $\goth{n}=\goth{m}B$ (o\`u $\goth{m}$,~$\goth{n}$ sont les
id\'eaux maximaux de~$A$,~$B$) la filtration $\goth{m}$-adique sur~$B$ est identique \`a sa filtration $\goth{n}$-adique et on aura
$$
\gr^*(B)=\gr^*(A)\otimes_k L.
$$
Il s'ensuit que $\gr^*(B)$ est une alg\`ebre de polyn\^omes sur~$L$ si et seulement si $\gr^*(A)$ est une alg\`ebre de polyn\^omes
sur~$k$. cqfd. (N.B.\ on n'a pas utilis\'e le fait que $L/k$ est
s\'eparable).

\begin{corollaire}
\label{I.9.2}
Soit $f\colon X\to Y$ un morphisme \'etale. Si $Y$ est
r\'egulier, $X$ l'est, la r\'eciproque \'etant vraie si $f$ est
surjectif.
\end{corollaire}

\setcounter{subsection}{1}

\begin{proposition}
\label{prop:I.9.2}
Soit $f\colon X\to Y$ un morphisme \'etale. Si $Y$ est r\'eduit,
il en est de m\^eme de~$X$, la r\'eciproque \'etant vraie si $f$
est surjectif.
\end{proposition}

Cela \'equivaut au
\begin{corollaire}
\label{I.9.3}
Soit $f\colon A\to B$ un homomorphisme local \'etale, $B$ \'etant
isomorphe \`a une
$A$-alg\`ebre localis\'ee d'une $A$-alg\`ebre de type fini.
Pour que $A$ soit r\'eduit, il faut et il suffit que $B$ le soit.
\end{corollaire}

La n\'ecessit\'e est triviale, puisque $A\to B$ est injectif ($B$
\'etant fid\`element plat sur~$A$). Suffisance: soient
$\goth{p}_i$ les id\'eaux premiers minimaux de~$A$, par
hypoth\`ese l'application naturelle $A\to\prod_i A/\goth{p}_i$ est
injective, donc tensorisant avec le $A$-module plat $B$, on trouve que
$B\to\prod_i B/\goth{p}_iB$ est injective, et on est ramen\'e \`a
prouver que les $B/\goth{p}_iB$ sont r\'eduits. Comme
$B/\goth{p}_iB$ est \'etale sur~$A/\goth{p}_i$,
on est ramen\'e au
cas $A$ int\`egre. Soit~$K$ son corps des fractions, alors $A\to K$
\'etant injectif, il en est de m\^eme ($B$~\'etant $A$-plat) de
$B\to B\otimes_A K$, on est ramener \`a prouver que ce dernier
anneau est r\'eduit. Or, $B$ \'etant localis\'ee d'une
$A$-alg\`ebre de type fini sur~$A$, est l'anneau local d'un point~$x$ d'un sch\'ema de type fini et \emph{étale} $X=\Spec(C)$ sur~$Y=\Spec(A)$, donc $B\otimes_A K$ est un anneau localis\'e (par
rapport \`a un ensemble multiplicativement stable convenable) de
l'anneau $C\otimes_A K$ de~$X\otimes_A K$. Comme $X\otimes_A K$ est
\'etale sur~$K$, son anneau est un produit fini de corps (extensions
s\'eparables de~$K$), il en est donc de m\^eme de~$B\otimes_A
K$, cqfd.

\begin{corollaire}
\label{I.9.4}
Soit $f\colon A\to B$ un homomorphisme local \'etale, supposons $A$
analytiquement r\'eduit
\index{analytiquement r\'eduit|hyperpage}(\ie le compl\'et\'e $\hat{A}$ de~$A$ sans \'el\'ements
nilpotents). Alors $B$ est analytiquement r\'eduit, et a~fortiori
r\'eduit.
\end{corollaire}

En effet,
$\hat{B}$
est \emph{fini} et \'etale sur~$\hat{A}$, et on applique~\ref{I.9.3}.

\begin{theoreme}
\label{I.9.5}
Soit $f\colon A\to B$ un homomorphisme local, $B$ \'etant isomorphe
\`a une alg\`ebre localis\'ee d'une $A$-alg\`ebre de type
fini. Alors:
\begin{enumerate}
\item[(i)] Si $f$ est \'etale, $A$ est normal si et seulement si $B$
l'est.
\item[(ii)] Si $A$ est normal, $f$ est \'etale si et seulement si
$f$ est injectif et net (et alors $B$ est normal par \textup{(i)}).
\end{enumerate}
\end{theoreme}

Nous allons donner deux d\'emonstrations diff\'erentes de (i), la
premi\`ere utilise certaines des propri\'et\'es des morphismes
plats quasi-finis (rappel\'es au d\'ebut du num\'ero) sans
utiliser~\ref{I.7.6} (et par l\`a, le Main Theorem); c'est l'inverse pour la
deuxi\`eme d\'emonstration. Enfin, pour~(ii) il semble qu'on ait
besoin du Main Theorem en tous cas.

\subsubsection*{Premi\`ere d\'emonstration}
On utilise la condition n\'ecessaire et suffisante suivante
de normalit\'e d'un anneau local noeth\'erien $A$ de dimension
$\neq0$.

\begin{enonce*}{Crit\`ere de Serre}
\index{Serre (crit\`ere de)|hyperpage}\textup{(i)}~Pour tout id\'eal premier $\goth{p}$ de~$A$ de rang~$1$,
$A_{\goth{p}}$ est normal (ou ce qui revient au m\^eme,
r\'egulier); \textup{(ii)}~Pour tout id\'eal premier $\goth{p}$ de~$A$ de
rang~$\geq 2$, on a profondeur $A_{\goth{p}}$ $\geq2$.\footnote{\Cf EGA~IV~5.8.6.}
\end{enonce*}

Nous admettrons ici ce crit\`ere, qui est cens\'e figurer au
paragraphe
des plats. Son principal avantage est qu'il ne suppose pas a priori
$A$ r\'eduit, ni a~fortiori int\`egre. Ici, on peut d\'ej\`a
supposer $\dim{A}=\dim{B}\neq0$.

D'apr\`es les rappels du d\'ebut du num\'ero, les id\'eaux
premiers $\goth{p}$ de~$A$ qui sont de rang~$1$ (\resp de
rang~$\ge2$) sont exactement les traces sur~$A$ des id\'eaux
premiers~$\goth{q}$ de~$B$ qui sont de rang~$1$ (\resp de
rang~$\ge2$). Enfin, si $\goth{p}$ et $\goth{q}$ se correspondent,
$B_{\goth{q}}$ est \'etale sur~$A_{\goth{p}}$, donc a m\^eme
profondeur que $A_{\goth{p}}$, et est r\'egulier si et seulement si
$A_{\goth{p}}$ l'est (\ref{I.9.1}). Appliquant le crit\`ere de Serre, on
trouve que $A$ est normal si et seulement si $B$ l'est.

\subsubsection*{Deuxi\`eme d\'emonstration}
Supposons $B$ normal, soit $L$ son corps des fractions, $K$ celui de
$A$ ($A$~est int\`egre puisque $B$ l'est). On a vu dans la
d\'emonstration de~\ref{I.9.3} que $B\otimes_A K$ est un compos\'e fini de
corps, comme il est contenu dans $L$ c'est un corps, et comme il
contient $B$ c'est $L$. Un \'el\'ement de~$K$ entier sur~$A$ est
entier sur~$B$, donc est dans $B$ puisque
$B$
est normal, donc dans $A$ car $B\cap K=A$ (comme il r\'esulte du
fait que $B$ est fid\`element plat sur~$A$).

Supposons maintenant $A$ normal, prouvons que $B$ l'est. En vertu
de~\ref{I.7.6} on aura $B=B^\prime_{\goth{n}}$, o\`u $B^\prime=A[t]/FA[t]$,
$F$ et $\goth{n}$ \'etant comme dans~\ref{I.7.6}. Donc $L=B\otimes_A K$
sera un localis\'e de~$B^\prime\otimes_A K=K[t]/FK[t]$, et un
produit de corps; extensions finies s\'eparables de~$K$ ce dernier
produit (comme chaque fois qu'on localise un anneau artinien, ici
$B^\prime_K$ par rapport \`a un ensemble multiplicativement stable)
est un facteur direct de~$B^\prime_K$, correspondant donc \`a une
d\'ecomposition $F=F_1F_2$ dans $K[t]$, le g\'en\'erateur de~$L$
correspondant \`a $t$ \'etant annul\'e d\'ej\`a par~$F_1$.
Or, \emph{$A$ étant normal, les $F_i$ sont dans $A[t]$} (supposant
qu'ils sont unitaires). Remarquant que $B\to L=B\otimes_A K$ est
injectif ($A\to K$ l'\'etant et $B$ \'etant plat sur~$A$) il
s'ensuit qu'on aura d\'ej\`a
$F_1(u)=0$, avec $u$ la
classe de~$t$ dans~$L$.
Supposant qu'on ait pris $F$ de degr\'e minimum, il s'ensuivra que
$F_2=1$ (N.B.\ on aura
$F^\prime(u)=F^\prime_1(u)F_2(u)+F_1(u)F^\prime_2(u)=F^\prime_1(u)F_2(u)$
puisque $F_1(u)=0$, d'o\`u $F^\prime_1(u)\neq0$ puisque
$F^\prime(u)\neq0$.)

Donc
on a
\begin{equation*}
\label{eq:I.9.5.*}
\tag{$*$} L=B\otimes_A K=K[t]/FK[t]
\end{equation*}
$F$ \'etant par suite un polyn\^ome s\'eparable dans $K[t]$
(mais \'evidemment par n\'ecessairement dans $A[t]$). (N.B.\ Pour
l'instant, on a seulement montr\'e, essentiellement, que dans~\ref{I.7.6} on
peut choisir $F$ et
$\goth{n}$
de telle fa\c con que --- avec les notations prises ici --- $B^\prime\to B^\prime_{\goth{n}}=B$
soit \emph{injectif}; on s'est servi pour cela de la normalit\'e de
$A$; je ne sais pas si cela reste vrai sans hypoth\`ese de
normalit\'e).

Rappelons maintenant le lemme bien connu, extrait du Cours de Serre de
l'an dernier:

\begin{lemme}
\label{I.9.6}
Soient $K$ un anneau, $F\in K[t]$ un polyn\^ome unitaire
s\'eparable, $L=K[t]/FK[t]$, $u$~la classe de~$t$ dans~$L$ (de sorte
que $F^\prime(u)$ est un \'el\'ement inversible de~$L$). Alors on
a les formules (o\`u $n=\deg{F}$):
\begin{enumerate}
\item[]$\trace_{L/K} u^i/F^\prime(u)=0$\ \ si\ \
$0\le i< n-1$,
\item[]$\trace_{L/K} u^{n-1}/F^\prime(u)=1$.
\end{enumerate}
\end{lemme}

\begin{corollaire}
\label{I.9.7}
Le d\'eterminant de la matrice $(u^j\cdot u^i/F^\prime(u))_{0\le i,j\le
n-1}$ est \'egal \`a
$(-1)^{n(n-1)/2}$,
donc inversible dans tout
sous-anneau $A$ de~$K$.
\end{corollaire}

\begin{corollaire}
\label{I.9.8}
Soient
$A$ un sous-anneau de~$K$, $V$ le $A$-module engendr\'e par les
$u^i$~($0\le i\le n-1$) dans $L$, $V^\prime$ le sous-$A$-module de~$L$
form\'e des $x\in L$ tels que $\trace_{L/K}(xy)\in A$ pour tout
$y\in V$ (\ie pour $y$ de la forme $u^i$, $0\le i\le n-1$). Alors
$V^\prime$ est le $A$-module ayant pour base les $u^i/F^\prime(u)$
($0\le i\le n-1$).
\end{corollaire}

\begin{corollaire}
\label{I.9.9}
Supposons que $K$ soit le corps des fractions d'un anneau int\`egre
normal $A$, $F$ ayant ses coefficients dans $A$. Alors avec les
notations de \ref{I.9.8}, $V^\prime$ contient la cl\^oture normale
$A^\prime$ de~$A$ dans $L$, qui est donc contenue dans
$A[u]/F^\prime(u)$ et a~fortiori dans $A[u][F^\prime(u)^{-1}]$.
\end{corollaire}

Appliquons ce dernier corollaire \`a la situation que nous avions
obtenue dans la d\'e\-mon\-stra\-tion: comme $F^\prime(u)$ est
inversible dans $B$ qui contient $A[u]$, $B$ contient $A^\prime$.
D'apr\`es le Main Theorem, (ou \`a partir du fait que
$B=A[u]_{\goth{n}}$) $B$ est une alg\`ebre localis\'ee de
$A^\prime$. Comme $A^\prime$ est normal, il en est de m\^eme
de~$B$.

\subsubsection*{D\'emonstration de \textup{(ii)}}
On
proc\`ede comme dans la d\'emonstration qui pr\'ec\`ede pour
prouver qu'on peut, dans~\ref{I.7.6}, choisir $F$ de telle fa\c con que l'on
ait encore ($*$). Le seul obstacle a priori est que, $B$ n'\'etant
plus suppos\'e plat sur~$A$, on ne peut plus affirmer que $B\to L$
est injectif, de sorte que le raisonnement ne s'appliquera a priori
qu'\`a l'image $B_1$ de~$B$ par ledit homomorphisme. Il s'ensuit
aussit\^ot que $B_1$ est plat sur~$A$ (comme localis\'ee d'une
alg\`ebre libre sur~$A$). En vertu de \ref{I.4.8} le morphisme $B\to B_1$
est \'etale, donc un isomorphisme, ce qui ach\`eve la
d\'emonstration.

(Du point de vue r\'edaction, il faudrait intervertir les deux
derni\`eres d\'emonstrations, et mettre dans un num\'ero \`a
part les calculs formels du lemme et de ses corollaires).

\begin{corollaire}
\label{I.9.10}
Soit $f\colon X\to Y$ un morphisme \'etale. Si $Y$ est normal, $X$
l'est, la r\'eciproque est vraie si $f$ est
surjectif.
\end{corollaire}

\begin{corollaire}
\label{I.9.11}
Soit $f\colon X\to Y$ un morphisme dominant, $Y$ \'etant normal et
$X$ connexe. Si $f$ est net, $f$ est \'etale, donc $X$ est normal
et par suite (\'etant connexe) irr\'eductible.
\end{corollaire}

Soit $U$ l'ensemble des points o\`u $f$ est \'etale, il est
ouvert, et il suffit de montrer qu'il est aussi ferm\'e et non vide.
$U$ contient l'image inverse du point g\'en\'erique de~$Y$ (car
pour une alg\`ebre sur un corps, non ramifi\'e $=$ \'etale) donc
($X$ dominant $Y$) est non vide. Si $x$ appartient \`a
l'adh\'erence de~$U$, alors il appartient \`a l'adh\'erence
d'une composante irr\'eductible $U_i$ de~$U$, donc \`a une
composante irr\'eductible $X_i\overset{v}{=}\bar{U}_i$
de~$X$ qui rencontre $U$, et par suite domine $Y$ (car toute
composante de~$U$, plat sur~$Y$, domine $Y$). Par suite, si $y$ est
la projection de~$x$ sur~$Y$, $\cal{O}_y\to\cal{O}_x$ est
\emph{injectif} (compte tenu que $\cal{O}_y$ est int\`egre). Comme
$\cal{O}_y$ est normal et $\cal{O}_y\to\cal{O}_x$ net, on conclut
\`a l'aide de \ref{I.9.5}(ii).

\begin{corollaire}
\label{I.9.12}
Soit $f\colon X\to Y$ un morphisme de type fini dominant, avec $Y$
normal et $X$ irr\'eductible. Alors l'ensemble des points o\`u
$f$ est \'etale est identique au compl\'ementaire du support de
$\Omega^1_{X/Y}$, \ie au compl\'ementaire du
sous-pr\'esch\'ema de~$X$ d\'efini par l'id\'eal
diff\'erente $\goth{d}_{X/Y}$.
\end{corollaire}

(C'est
cela l'\'enonc\'e \og moins trivial\fg auquel il \'etait fait
allusion dans la remarque du~\No\ref{I.4}.)

\begin{remarquestar}
On se gardera de croire qu'un rev\^etement \'etale connexe d'un
sch\'ema irr\'eductible soit lui-m\^eme irr\'eductible, quand
on ne suppose pas la base normale. Cette question sera
\'etudi\'ee au~\No \ref{I.11}.
\end{remarquestar}

\section{Rev\^etements \'etales d'un sch\'ema normal}
\label{I.10}

\begin{proposition}
\label{I.10.1}
Soit $X$ un pr\'esch\'ema \'etale s\'epar\'e sur~$Y$ normal
connexe de corps~$K$. Alors les composantes connexes

$X_i$ de~$X$ sont int\`egres, leurs corps $K_i$ sont des extensions
finies s\'eparables de~$K$, $X_i$ s'identifie \`a une partie
ouverte non vide du normalis\'e de~$X$ dans $K_i$ (donc $X$ \`a
une partie ouverte dense du normalis\'e de~$Y$ dans $R(X)=L=\prod
K_i$).
\end{proposition}

D'apr\`es \ref{I.9.10} $X$ est normal, a~fortiori ses anneaux locaux sont
int\`egres, donc les composantes connexes de~$X$ sont
irr\'eductibles. Comme $X_i$ est normal, et fini et dominant
au-dessus de~$Y$, il r\'esulte d'un cas particulier (\`a peu
pr\`es trivial d'ailleurs) du Main Theorem que $X_i$ est un ouvert
du normalis\'e de~$X$ dans le corps $K_i$ de~$X$.

\begin{corollaire}
\label{I.10.2}
Sous les conditions \ref{I.10.1}, $X$ est fini sur~$Y$ (\ie un
rev\^etement \'etale de~$Y$) si et seulement si $X$ est isomorphe
au normalis\'e $Y'$ de~$Y$ dans $L=R(X)$ (anneau des fonctions
rationnelles sur~$X$).
\end{corollaire}

En effet, on sait que ce normalis\'e est fini sur~$Y$ ($Y$ \'etant
normal et $R/K$ s\'eparable), inversement si $X$ est fini sur~$Y$ il
l'est sur~$Y'$, donc son image dans $Y'$ est ferm\'ee, d'autre part
elle est dense.

Une alg\`ebre $L$ de rang fini sur~$K$ sera dite \emph{non
ramifiée sur~$X$} (ou simplement non ramifi\'ee sur~$K$, si $X$
est sous-entendu) si $L$ est une alg\`ebre s\'eparable sur~$K$,
\ie compos\'ee directe d'extensions s\'eparables $K_i$, et si le
normalis\'e $Y'$ de~$Y$ dans $L$ (somme disjointe des normalis\'es
de~$Y$ dans les $K_i$) est non ramifi\'e ($=$ \'etale par \ref{I.9.11})
sur~$Y$. Donc:

\begin{corollaire}
\label{I.10.3}
Pour tout $X$ fini sur~$Y$ et dont toute composante irr\'eductible
domine $Y$, soit $R(X)$ l'anneau des fonctions rationnelles sur~$X$
(produit des anneaux
locaux des points g\'en\'eriques des composantes irr\'eductibles
de~$X$), de sorte que $X\mto R(X)$ est un foncteur, \`a valeurs
dans les alg\`ebres de rang fini sur~$K=R(Y)$. Ce foncteur
\'etablit une \'equivalence de la cat\'egorie des rev\^etement
\'etales connexes de~$Y$ avec la cat\'egorie des extensions $L$ de
$K$ non ramifi\'ees sur~$Y$.
\end{corollaire}

Le foncteur inverse est le foncteur normalisation.

Supposons $Y$ affine, donc d\'efini par un anneau normal $A$ de corps
des fractions~$K$. Soit $L$ une extension finie de~$K$ compos\'e
directe de corps, alors par d\'efinition la normalis\'ee $Y'$ de
$Y$ dans $L$ est isomorphe \`a $\Spec(A')$, o\`u $A'$ est le
normalis\'e de~$A$ dans~$L$. Dire que $L$ est non ramifi\'e sur~$Y$ signifie que $A'$ est non ramifi\'e (ou encore: \'etale)
sur~$A$.
Si
$A$ est local, il revient au m\^eme de dire que les anneaux locaux
$A^\prime_{\goth{n}}$ (o\`u
$\goth{n}$
parcourt l'ensemble fini des id\'eaux maximaux de~$A'$, \ie de ses
id\'eaux premiers induisant l'id\'eal maximal $\goth{m}$ de~$A$)
soient non ramifi\'es ($=$ \'etales) sur

l'anneau local~$A$. Enfin, notons aussi que le crit\`ere par le
discriminant (\ref{I.4.10}) peut aussi s'appliquer dans cette situation (plus
g\'en\'eralement, une variante dudit crit\`ere devrait
s'\'enoncer ainsi, sans condition pr\'eliminaire de platitude
lorsque $X$ domine $Y$, $Y$ \'etant n\'eanmoins suppos\'e
localement int\`egre: $A\to B$ et $B\to B\otimes_A K =L$ sont
injectifs --- alors $\trace_{L/K}$ est d\'efinie --- et
$\trace_{L/K}(xy)$ induit une \emph{forme bilinéaire fondamentale}
$B\times B\to A$, \ie il existe des $x_i\in B$ ($1\le i\le n$, $n=$
rang de~$L$ sur~$K$) tels que $\trace(x_ix_j)\in A$ pour tout $i$,~$j$
et
$\det(\trace(x_ix_j))$
est inversible dans~$A$).

Le sorite (4.6) implique aussit\^ot le sorite de la non ramification
dans le cadre classique:

\begin{proposition}
\label{I.10.4}
Soit
$Y$ un pr\'esch\'ema normal int\`egre, de corps~$K$.
\textup{(i)} $K$ est non ramifi\'e
sur~$Y$.
\textup{(ii)} Si $L$ est une extension de~$K$ non ramifi\'ee sur~$Y$, si $Y'$
est un pr\'esch\'ema normal de corps $L$ et dominant $Y$ (par
exemple le normalis\'e de~$Y$ dans~$L$) et $M$ une extension de~$L$
non ramifi\'ee sur~$Y'$, alors $M/K$ est non ramifi\'ee sur~$X$
(\emph{transitivité} de la non ramification). \textup{(iii)} Soit $Y'$ un
pr\'esch\'ema normal int\`egre dominant $Y$, de corps $K'/K$; si
$L$ est une extension de~$K$ non ramifi\'ee sur~$Y$, alors
$L\otimes_K K'$ est une extension de~$K'$ non ramifi\'ee sur~$Y'$
(propri\'et\'e de \emph{translation}).
\end{proposition}

De
plus:

\begin{corollaire}
\label{I.10.5}
Sous les conditions de \textup{(iii)}, si $Y=\Spec(A)$, $Y'=\Spec(A')$, alors
le normalis\'e
$\bar{A'}$
de~$Y'$ dans $L'=L\otimes_K K'$ s'identifie
\`a $\bar{A}\otimes_A A'$, o\`u $\bar{A}$ est le normalis\'e de
$A$ dans $L$.
\end{corollaire}

Habituellement, les gens (qui r\'epugnent \`a la consid\'eration
d'anneaux non int\`egres, fussent-ils compos\'es directs de corps)
\'enoncent la propri\'et\'e de translation sous la forme (plus
faible) suivante:

\begin{corollaire}
\label{I.10.6}
Sous les conditions de \textup{(iii)}, soit $L_1$ une \emph{extension
composée} de~$L/K$ (non ramifi\'ee sur~$Y$) et de~$K'/K$. Alors
$L_1/K'$ est non ramifi\'ee sur~$Y'$. Dans le cas o\`u
$Y=\Spec(A)$, $Y'=\Spec(A')$, on aura de plus
$$
\overline{A'}=A[\overline{A},A']
$$
\ie l'anneau $\overline{A'}$ normalis\'e de~$A'$ dans $L_1$ est la
$A$-alg\`ebre engendr\'ee par $A'$ et le normalis\'e
$\overline{A}$ de~$A$ dans $L$.
\end{corollaire}

Ce dernier fait est d'ailleurs faux sans hypoth\`ese de non
ramification, m\^eme dans le cas d'extensions compos\'ees de corps
de nombres...

Pour terminer ce num\'ero, nous allons donner l'interpr\'etation
de la notion de rev\^etement \'etale correspondant \`a l'image
intuitive de cette notion: il doit y avoir le \og nombre maximum\fg de
points au-dessus du point consid\'er\'e $y\in Y$, et en
particulier il ne doit pas y avoir \og plusieurs points confondus\fg
au-dessus de~$y$. Pour d\'emontrer les r\'esultats dans ce sens
avec toute la g\'en\'eralit\'e d\'esirable, nous allons
admettre ici la proposition~\ref{I.10.7} plus bas (dont la d\'emonstration
sera dans le multiplodoque, Chap.\, IV, par.\, 15, et utilise la technique
des ensembles constructibles de Chevalley, et un petit peu de
th\'eorie de descente...).

Un morphisme de type fini $f\colon X\to Y$ est dit
\emph{universellement ouvert} si pour toute extension de la base
$Y'\to Y$ (avec $Y'$ localement noeth\'erien) le morphisme $f'\colon
X'=X\times_Y Y'\to Y'$ est ouvert, \ie transforme ouverts en
ouverts. On peut d'ailleurs se borner au cas o\`u $Y'$ est de type
fini sur~$Y$ (et m\^eme o\`u $Y'$ est de la forme
$Y[t_1,\dots,t_r]$, o\`u les $t_i$ sont des ind\'etermin\'ees).
Un morphisme universellement ouvert est a~fortiori ouvert (la
r\'eciproque \'etant fausse), d'autre part si $f$ est ouvert, $X$
et $Y$ \'etant irr\'eductibles, alors toutes les composantes de
toutes les fibres de~$f$ ont m\^eme dimension (savoir la dimension de
la fibre g\'en\'erique
$f^{-1}(z)$, $z$ le point g\'en\'erique
de~$Y$). Enfin si $Y$ est normal, cette derni\`ere condition
implique d\'ej\`a que $f$ est \emph{universellement} ouvert
(th\'eor\`eme de Chevalley). Il s'ensuit par exemple que si
$f\colon X\to Y$ est un morphisme \emph{quasi-fini}, avec $Y$ normal
irr\'eductible, alors $f$ est universellement
ouvert
(ou encore:
ouvert)
si et seulement si toute composante irr\'eductible de~$X$
domine~$Y$. Rappelons aussi qu'un morphisme plat (de type fini)
\'etant ouvert, est aussi universellement ouvert. Ces
pr\'eliminaires pos\'es, \og rappelons\fg la

\begin{proposition}
\label{I.10.7}
Soit $f\colon X\to Y$ un morphisme quasi-fini s\'epar\'e
universellement ouvert. Pour tout $y\in Y$, soit $n(y)$ le \og nombre
g\'eom\'etrique de points de la fibre $f^{-1}(y)$\fg, \'egal
\`a la somme des degr\'es s\'eparables des extensions
r\'esiduelles
$\kres(x)/\kres(y)$, pour les points $x\in f^{-1}(y)$. Alors
la fonction $y\mto n(y)$ sur~$Y$ est semi-continue
sup\'erieurement. Pour qu'elle soit constante au voisinage du point
$y$ (\ie pour qu'on ait $n(y)=n(z_i)$, o\`u les $z_i$ sont les
points g\'en\'eriques des composantes irr\'eductibles de~$Y$ qui
contiennent $y$) il faut qu'il existe un voisinage $U$ de~$y$ tel que
$X|U$ soit \emph{fini} sur~$U$.\footnote{\Cf EGA IV 15.5.1}
\end{proposition}

\begin{corollaire}
\label{I.10.8}
Si $y\mto n(y)$ est constante et $Y$ g\'eom\'etriquement
unibranche\footnote{Pour la d\'efinition, \cf ci-dessous
\no \ref{I.11}, p.\, \pageref{I.11}}, les composantes irr\'eductibles
de~$X$ sont disjointes.
\end{corollaire}

\begin{proposition}
\label{I.10.9}
Soit $f\colon X\to Y$ un morphisme \emph{étale} s\'epar\'e.
Avec les notations~\ref{I.10.7} la fonction $y\mto n(y)$ est
semi-continue sup\'erieurement. Pour qu'elle soit constante au
voisinage du point $y$, (\ie pour qu'on ait $n(y)=n(z_i)$, o\`u
les $z_i$ sont les points g\'en\'eriques des composantes
irr\'eductibles de~$Y$ qui contiennent $y$) il faut et il suffit
qu'il existe un voisinage ouvert $U$ de~$y$ tel que $X|U$ soit fini
sur U, \ie soit un \emph{revêtement étale} de~$U$.
\end{proposition}

\begin{corollaire}
\label{I.10.10}
Pour qu'un morphisme \'etale s\'epar\'e $f\colon X\to Y$, $Y$
connexe, soit fini (\ie fasse de~$X$ un \emph{revêtement
étale} de~$Y$) il faut et il suffit que toutes les fibres de~$f$
aient m\^eme nombre g\'eom\'etrique de points.
\end{corollaire}

Dans~\ref{I.10.7}
et son corollaire, il n'y avait pas d'hypoth\`ese de normalit\'e
sur~$Y$. Si on fait une telle hypoth\`ese, on trouve
l'\'enonc\'e plus fort (pris le plus souvent comme
d\'efinition de la non ramification d'un rev\^etement):

\begin{theoreme}
\label{I.10.11}
Soit
$f\colon X\to Y$ un morphisme \emph{quasi-fini} s\'epar\'e.
On suppose que~$Y$ est \emph{irréductible}, que toute composante
de~$X$ domine $Y$, que $X$ soit r\'eduit (\ie $\cal{O}_X$ sans
\'el\'ements nilpotents). Soit $n$ le degr\'e de~$X$ sur~$Y$
(somme des degr\'es, sur le corps $K$ de~$Y$, des corps $K_i$ des
composantes irr\'eductibles $X_i$ de~$X$). Soit $y$ un point normal
de~$Y$. Alors le nombre g\'eom\'etrique $n(y)$ de points de~$X$
au-dessus de~$y$ est $\le n$, l'\'egalit\'e ayant lieu si et
seulement si il existe un voisinage ouvert $U$ de~$y$ tel que $X|U$
soit un \emph{revêtement étale} de~$U$.
\end{theoreme}

Le \og seulement si\fg \'etant trivial, prouvons le \og si\fg. Soit $z$ le
point g\'en\'erique de~$Y$, on a $n(z)=$ (somme des degr\'es
s\'eparables des $K_i/K$) $\le n$ et par \ref{I.10.7} on a $n(y)\le n(z)\le
n$, l'\'egalit\'e impliquant que $X|U$ est \emph{fini} sur~$U$
pour un voisinage~$U$ convenable de~$y$. On peut donc supposer $X$
fini sur~$Y$ et la fonction $n(y')$ sur~$Y$ constante. Enfin par~\ref{I.10.8}
$X$ est alors r\'eunion disjointe de ses composantes
irr\'eductibles et pour prouver qu'il est non ramifi\'e en $y$, on
est ramen\'e au cas o\`u $X$ est irr\'eductible, donc
int\`egre. Enfin on peut supposer $Y=\Spec(\cal{O}_y)$. Le
th\'eor\`eme se r\'eduit alors \`a l'\'enonc\'e classique
suivant:

\begin{corollaire}
\label{I.10.12}
Soient $A$ un anneau local normal (noeth\'erien comme toujours) de
corps $K$, $L$ une extension finie de~$K$ de degr\'e $n$, degr\'e
s\'eparable $n_s$, $B$ un sous-anneau de~$L$ fini sur~$A$, de corps
des fractions
$L$,
$\goth{m}$ l'id\'eal maximal de~$A$ et $n'$ le
degr\'e s\'eparable de~$B/\goth{m}B$ sur~$A/\goth{m}A=k$ ($=$
somme des degr\'es s\'eparables des extensions r\'esiduelles de
cet anneau). On a $n'\le n_s$ et a~fortiori $n'\le n$. Cette
derni\`ere in\'egalit\'e est une \'egalit\'e si et seulement
si $B$ est non~ramifi\'e ($=$ \'etale) sur~$A$.
\end{corollaire}

Il reste seulement \`a montrer que $n'=n$ implique que $B$ est
\'etale sur~$A$. Rappelons la d\'emonstration quand $k$ est
infini: on doit seulement montrer que $R=B/\goth{m}B$ est
s\'eparable sur~$k$; s'il n'en \'etait pas ainsi il en
r\'esulterait (par un lemme connu) qu'il existe un \'el\'ement
$a$ de~$R$ dont le polyn\^ome minimal sur~$k$ est de degr\'e
$>n'$. Cet \'el\'ement provient d'un \'el\'ement $x$ de~$B$,
dont le polyn\^ome minimal sur~$K$ (en tant qu'\'el\'ement de
$L$) est de degr\'e $\le n$; d'autre part ce dernier a ses
coefficients dans~$A$ puisque $A$ est normal, et donne donc par
r\'eduction mod $\goth{m}$ un polyn\^ome unitaire $F\in k[t]$, de
degr\'e $\le n=n'$, tel que $F(a)=0$, absurde.

Dans
le cas g\'en\'eral ($k$ pouvant \^etre fini), reprenant le
langage g\'eom\'etrique, on consid\`ere $Y'=\Spec(A[t])$ qui est
fid\`element plat sur~$Y$, et le point g\'en\'erique $y'$ de la
fibre $\Spec(k[t])$ de~$Y'$ sur~$y$. Alors $X$ est net sur~$Y$ en $y$
si et seulement si $X'=X\times_Y Y'=\Spec(B[t])$ est net en $y'$ sur~$Y'$, comme on constate aussit\^ot. D'autre part, d'apr\`es le
choix de~$y'$, son corps r\'esiduel est $k(t)$ donc infini. Comme
$y'$ est un point normal de~$Y'$, on est ramen\'e au cas
pr\'ec\'edent.







\section{Quelques compl\'ements}
\label{I.11}

Nous avons d\'ej\`a dit qu'un rev\^etement \'etale connexe
d'un sch\'ema int\`egre n'est pas n\'e\-ces\-sai\-rement
int\`egre. Voici deux exemples de ce fait.

a) Soit $C$ une courbe alg\'ebrique \`a point double
ordinaire~$x$, $C^\prime$ sa normalis\'ee, $a$ et~$b$ les deux
points de~$C^\prime$ au-dessus de~$x$. Soient $C_i^\prime$ ($i=1$,
$2$) deux copies de~$C^\prime$, $a_i$ et~$b_i$ le point de
$C_i^\prime$ qui correspond \`a $a$ \resp $b$. Dans la courbe somme
$C_1^\prime \amalg C_2^\prime$, identifions
$a_1$
et~$b_2$ d'une part, $a_2$ et~$b_1$ d'autre part (on laisse au lecteur
le soin de pr\'eciser le processus d'identification; il sera
expliqu\'e au Chap.\, VI du multiplodoque mais, dans le cas des
courbes sur un corps alg\'ebriquement clos, est trait\'e dans le
livre de Serre sur les courbes alg\'ebriques). On trouve une courbe
$C^{\prime\prime}$ \emph{connexe} et \emph{réductible}, qui est un
rev\^etement \'etale de degr\'e~$2$ de~$C$. Le lecteur
v\'erifiera que de fa\c con g\'en\'erale, les rev\^etements
\'etales connexes \og galoisiens\fg $C^{\prime\prime}$ de~$C$ dont
l'image inverse $C^{\prime\prime} \times_C C^\prime$ est un
rev\^etement \emph{trivial} de~$C^\prime$ (\ie isomorphe \`a la
somme d'un certain nombre de copies de~$C^\prime$) sont \og cycliques\fg
de degr\'e~$n$, et pour tout entier~$n>0$, on peut construire un
rev\^etement \'etale connexe cyclique de degr\'e~$n$. Dans le
langage du groupe fondamental qui sera d\'evelopp\'e plus tard,
cela signifie que le quotient de~$\pi_1(C)$ par le sous-groupe
invariant ferm\'e engendr\'e par l'image de~$\pi_1(C^\prime) \to
\pi_1(C)$ (homomorphisme induit par la projection) est isomorphe au
compactifi\'e de~$\ZZ$. De fa\c con plus pr\'ecise, on doit
pouvoir montrer que le groupe fondamental de~$C$ est isomorphe au
produit libre (topologique) du groupe fondamental
de~$C^\prime$
par le compactifi\'e de~$\ZZ$. Notons que ce sont des questions
de ce genre qui ont donn\'e naissance \`a la \og th\'eorie de la
descente\fg pour les sch\'emas.

b) Soit $A$ un anneau local complet int\`egre, on sait que son
normalis\'e $A^\prime$ est fini
sur~$A$ (Nagata), donc c'est un anneau semi-local complet, donc local
puisqu'il est int\`egre. Supposons que l'extension r\'esiduelle
$L/k$ qu'il d\'efinit soit non radicielle (dans le cas contraire, on
dira que $A$ est g\'eom\'etriquement unibranche, \cf plus bas). Ce
sera le cas par exemple pour l'anneau
$\RR[[s,t]]/(s^2+t^2)\RR[[s,t]]$,
o\`u $\RR$ est le corps
des r\'eels. Soit alors $k^\prime$ une extension galoisienne finie
de~$k$ telle que $L \otimes_k k^\prime$ se d\'ecompose; et soit~$B$
une alg\`ebre finie \'etale sur~$A$ correspondant \`a
l'extension r\'esiduelle~$k^\prime$ (rappelons que $B$ est
essentiellement unique). Alors $B^\prime = A^\prime \otimes_A B$
sur~$B$ a l'alg\`ebre r\'esiduelle $L \otimes_k k^\prime$ qui
n'est pas locale, donc $B^\prime$ n'est pas un anneau local donc
(\'etant complet) a des diviseurs de~$0$. Comme il est contenu dans
l'anneau total des fractions de~$B$ (car libre sur~$A^\prime$ donc
sans torsion sur~$A^\prime$ donc sans torsion sur~$A$, donc contenu
dans $B^\prime \otimes_A K = B^\prime_{(K)} = A^\prime_{(K)} \otimes_K
B_{(K)} = B_{(K)}$ puisque $A^\prime_{(K)} = K$) il s'ensuit que $B$
n'est pas int\`egre. Dans le cas de l'anneau
$\RR[s,t]/(s^2+t^2)\RR[s,t]$, prenant $k^\prime/k=
\CC/\RR$, on trouve pour~$B$ l'anneau local de deux droites
s\'ecantes du plan en leur point d'intersection.

Notons d'ailleurs que s'il existe un rev\^etement connexe \'etale
$X$ de~$Y$ int\`egre qui ne soit pas irr\'eductible, alors toute
composante irr\'eductible de~$X$ donne un exemple d'un
rev\^etement non ramifi\'e $X^\prime$ de~$Y$, dominant~$Y$, qui
n'est pas \'etale sur~$Y$. Dans le cas de l'exemple~a), on obtient
ainsi que $C^\prime$ est non ramifi\'e sur~$C$, sans \^etre
\'etale en les deux points $a$ et~$b$ (comme on constate d'ailleurs
directement par inspection des compl\'et\'es des anneaux locaux
de~$x$ et~$a$: du point de vue \og formel\fg, $C^\prime$ au point~$a$
s'identifie \`a un sous-sch\'ema ferm\'e de~$C$ au point~$x$,
savoir l'une des deux \og branches\fg de~$C$ passant par~$x$).

Dans a) et~b), on voit que la non validit\'e des conclusions
de~\ref{I.9.5}.\, (i) et~(ii) est li\'ee directement au fait qu'un
point de~$Y$ \og \'eclate\fg en des points \emph{distincts} du
normalis\'e (dans~b),
le fait que l'extension r\'esiduelle soit non radicielle
doit \^etre interpr\'et\'ee g\'eom\'etriquement de cette
fa\c con). De fa\c con pr\'ecise, nous dirons qu'un anneau local
int\`egre~$A$ est \emph{géométriquement unibranche}
\index{geometriquement unibranche@g\'eom\'etriquement unibranche|hyperpage}\index{unibranche (geometriquement)@unibranche (g\'eom\'etriquement)|hyperpage}si son normalis\'e n'a qu'un seul id\'eal maximal, l'extension
r\'esiduelle correspondante \'etant radicielle; un point~$y$ d'un
pr\'esch\'ema int\`egre est dit g\'eom\'etriquement
unibranche si son anneau local l'est. Exemples: un point normal, un
point de rebroussement ordinaire d'une courbe, etc. Il semble que si
$Y$ admet un point qui
n'est pas unibranche, il existe toujours
un
rev\^etement \'etale
connexe non irr\'eductible de~$Y$; c'est du moins ce que nous avons
montr\'e dans le cas~b), lorsque $Y$ est le spectre d'un anneau
local complet. On peut montrer par contre que \emph{si tous les points
de~$Y$ sont géométriquement unibranches, alors tout
$Y$-préschéma non ramifié connexe dominant~$Y$ est
étale} et irr\'eductible. La d\'emonstration reprend celle
de~\ref{I.9.5}, en utilisant la g\'en\'eralisation suivante du
th\'eor\`eme~\ref{I.8.3}, qui sera d\'emontr\'ee plus tard
\`a l'aide de la technique de
descente\footnote{\Cf IX~\ref{IX.4.10}. Pour une d\'emonstration
plus directe, \cf EGA~IV 18.10.3, utilisant une variante
de~\ref{I.9.5} pour des anneaux locaux g\'eom\'etriquement
unibranches.}:

\emph{Soit $Y^\prime \to Y$ un morphisme fini, radiciel, surjectif
(\ie ce qu'on pourrait appeler un \og homéomorphisme
universel\fg). Considérons le foncteur $X \mto X \times_Y
Y^\prime = X^\prime$ des $Y$-préschémas dans les
$Y^\prime$-préschémas. Ce foncteur induit une équivalence
de la catégorie des $Y$-schémas étales avec la
catégorie des $Y^\prime$-schémas étales.} On pourra
appliquer par exemple ce r\'esultat dans le cas o\`u $Y^\prime$
est le normalis\'e de~$Y$, $Y$ \'etant suppos\'e unibranche (et
$Y^\prime$ fini sur~$Y$, ce qui est vrai dans tous les cas qu'on
rencontre en pratique), ou au cas d'un $Y^{\prime\prime}$ \og en
sandwich\fg entre~$Y$ et son normalis\'e (qui n'a plus besoin
d'\^etre fini sur~$Y$).

